%% CV - ExxonMobil Geoscientist (Kuala Lumpur)
%% Tailored from cv/main_example.tex. Facts from 01-candidate-profile.md
%% (CV 2026, Skills PDF, publications list, LinkedIn).
%% Compile with: lualatex main_exxonmobil_geoscientist.tex

\documentclass[11pt,a4paper,sans]{moderncv}
\moderncvstyle{banking}
\moderncvcolor{blue}

\renewcommand*{\namefont}{\fontsize{34}{36}\bfseries\upshape}
\colorlet{firstnamecolor}{color1}
\colorlet{lastnamecolor}{color1}
\colorlet{namecolor}{color1}
\renewcommand*{\sectionstyle}[1]{{\sectionfont\color{color1}#1}}

\usepackage[utf8]{inputenc}
\AtEndPreamble{\hypersetup{
    colorlinks=true,
    linkcolor=blue,
    filecolor=magenta,
    urlcolor=blue,
    pdftitle={Nicholas Fernie - CV - ExxonMobil Geoscientist},
    pdfpagemode=UseNone,
}}
\usepackage[scale=0.77]{geometry}
\usepackage{import}
\usepackage{needspace}

\name{Nicholas}{Fernie}
\address{Teneriffe, Brisbane, Australia (open to relocating to Kuala Lumpur)}{}{}
\phone[mobile]{+61 421 776 104}
\email{fernienicholas@gmail.com}
\extrainfo{\href{https://www.linkedin.com/in/nic-fernie-28125084}{LinkedIn}, \href{https://github.com/NFernie}{GitHub}}

\begin{document}

\makecvtitle

\vspace{6pt}
\small{Development geologist with an MPhil in Geoscience, First Class Honours, and almost eight years at Santos supporting oil and tight-gas producing assets. Integrates reservoir characterization, sequence stratigraphy, and Petrel static reservoir models (uncertainty analysis, seismic depth conversion, static-to-dynamic handoff) to inform field development planning, opportunity generation, and drill-well programmes. Geosteered seven horizontal wells in four Cooper Oil fields and led a 22-well tight gas appraisal campaign.}

\section{Core Competencies}
\vspace{1pt}
\begin{itemize}

\item \textbf{Reservoir characterization / well log interpretation}: core reviews, well logs, sequence stratigraphy, and production geology (fault-driven and stratigraphic compartmentalization) applied to recovery-factor and field-screening decisions.

\item \textbf{Static reservoir models / uncertainty analysis}: Petrel field- and regional-scale models, workflow automation, seismic depth conversions, and static-to-dynamic handoff (tNavigator). Quantitative seismic interpretation is outside current practice.

\item \textbf{Field development / drill-well / opportunity generation}: factory-style tight-gas campaigns of 20+ wells including fracture stimulation; horizontal oil and gas well planning; conventional wells; 100+ onshore wells overseen.

\item \textbf{Structural analysis and subsurface mapping}: regional and field-level modelling, deposition-controlled compartmentalization, and geology-driven static models that honour core-to-analogue reservoir architecture.

\item \textbf{Integration with reservoir engineering}: production and dynamic data through static-to-dynamic workflows in support of depletion planning and reservoir management, not as a reservoir-engineer title.

\end{itemize}

\section{Professional Experience}
\vspace{3pt}
\begin{itemize}

\needspace{5\baselineskip}
\item{\cventry{Feb 2021 - Present}{Development Geologist}{Santos Ltd}{Brisbane, Australia}{}{\vspace{1pt}
\begin{itemize}
    \item Cooper Oil (2023-present): built geology-driven static reservoir models and geosteered 7 horizontal wells in four fields, targeting shallow marine and onshore fluvial reservoirs; honoured reservoir architecture from core-scale information to modern analogues.
    \item Cooper Central tight gas (from 2021): as project lead, delivered a 22-well tight gas appraisal campaign in Moomba; as non-lead geologist, identified conventional-type secondary targets to support the primary tight-gas objective.
    \item Opportunity generation and well planning across factory-style tight-gas programmes (20+ wells per campaign, including frac), plus horizontal oil/gas and conventional wells; overseen 100+ onshore wells.
    \item Sequence stratigraphy, full-hole core reviews, and compartmentalization analysis (fault and stratigraphic) to support field development planning and production optimization.
    \item Petrel uncertainty analysis, seismic depth conversions, regional and field-level modelling, and static-to-dynamic handoff.
\end{itemize}}}

\vspace{3pt}

\needspace{5\baselineskip}
\item{\cventry{Feb 2020 - Feb 2021}{Geomechanist}{Santos Ltd}{Brisbane, Australia}{}{\vspace{1pt}
\begin{itemize}
    \item Developed mechanical earth models (MEMs); completed offset-well reviews and frac reviews in support of well planning, horizontal drilling, and post-execution evaluation.
\end{itemize}}}

\vspace{3pt}

\needspace{5\baselineskip}
\item{\cventry{Jan 2019 - Jan 2020}{Graduate Wellsite Geologist}{Santos Ltd}{Adelaide, Australia}{}{\vspace{1pt}
\begin{itemize}
    \item Onshore and offshore wellsite geology: mudlogging, coring, wireline logging, and geological reporting during drilling.
\end{itemize}}}

\vspace{3pt}

\needspace{5\baselineskip}
\item{\cventry{Nov 2015 - Feb 2016}{Vacation Student}{Horizon Oil}{Sydney, Australia}{}{\vspace{1pt}
\begin{itemize}
    \item South China Sea oil and gas fields: wireline-log interpretation, seismic interpretation, core analysis, PVT data, structural mapping, and reservoir data with the subsurface team.
\end{itemize}}}

\end{itemize}

\section{Education}
\vspace{1pt}
\begin{itemize}

\needspace{5\baselineskip}
\item{\cventry{2017-2019}{Master of Philosophy, Geoscience}{University of Adelaide}{Adelaide, Australia}{}{\vspace{1pt}
Thesis: ``Thermal History of Central Australia: Cooper Basin, South Australia \& Anmatjira Ranges, Northern Territory: Insights from Apatite Fission Track \& U-Pb Thermochronology.''
}}

\vspace{3pt}

\needspace{5\baselineskip}
\item{\cventry{2013-2016}{BSc Geophysics and Applied Geology, First Class Honours}{University of Adelaide}{Adelaide, Australia}{}{\vspace{1pt}
Honours GPA 7.00. Honours thesis: low-temperature thermochronology of the Ghanaian margin and Equatorial Atlantic rifting. April 2018: ExxonMobil graduate field trip, Cape Liptrap, Victoria.
}}

\end{itemize}

\section{Selected Training}
\vspace{1pt}
\begin{itemize}
\item Seismic Geomorphology and Seismic Stratigraphy (2025); Reservoir Model Design (2025); Geostatistical Modelling for Reservoir Characterisation (2024); Development Planning for Mature Fields (2022); Non-marine Reservoirs \& Sequence Stratigraphy (2021).
\end{itemize}

\section{Languages}
\vspace{1pt}
\begin{itemize}
\item English (native). Spanish (intermediate).
\end{itemize}

\section{Publications}
\vspace{1pt}
\begin{itemize}
\item Nixon, A., Fernie, N., Glorie, S., Hand, M., Bendell, B. (2024). Thermal evolution and sediment provenance of the Cooper-Eromanga Basin: Insights from detrital apatite. \textit{Basin Research}. \href{https://doi.org/10.1111/bre.12843}{DOI}
\item Nixon, A., Glorie, S., Hasterok, D., Collins, A.~S., Fernie, N., Fraser, G. (2022). Low-temperature thermal history of the McArthur Basin: Influence of the Cambrian Kalkarindji Large Igneous Province on hydrocarbon maturation. \textit{Basin Research}. \href{https://doi.org/10.1111/bre.12691}{DOI}
\item Fernie, N., Glorie, S., Jessell, M., Collins, A.~S. (2018). Thermochronological insights into reactivation of a continental shear zone in response to Equatorial Atlantic rifting (northern Ghana). \textit{Scientific Reports}. \href{https://doi.org/10.1038/s41598-018-34769-x}{DOI}
\end{itemize}

\section{References}
\vspace{1pt}
\begin{itemize}
\item{Available upon request.}
\end{itemize}

\end{document}
